% ===================================================================
%  QUADRO N.º 6 — A Casa por dentro, os Móveis, a Comida, a Iluminação.
%
%  Ceci n'est PAS une transcription : c'est une traduction, faite en
%  2026, en portugais d'aujourd'hui. Voir texto/pt/10-quadro-01.tex
%  pour la consigne, qui vaut pour les seize fichiers.
% ===================================================================

%%K t06-apar-1 apar
\VUsaut{6.0mm}
\VUtitre{15.0pt}{60}{QUADRO N.º 6}
\VUsaut{1.4mm}
\VUtitre{11.4pt}{0}{A Casa por dentro, os Móveis, a Comida,}
\VUsaut{0.4mm}
\VUtitre{11.4pt}{0}{a Iluminação.}
\VUsaut{2.2mm}

%%K t06-apar-2 apar
A casa está terminada. O tio de João instalou-se na sua nova morada,
cuja distribuição já conhecemos. Vejamos agora como a mobilou.
Visitaremos primeiro:

%%K t06-tit-1 sub
\VUpk{10.2pt}{40}{O Quarto de dormir.}
\VUsaut{3.0mm}

%%K t06-c1-01-1 p
1. --- Chegamos a má hora, porque a criada está ocupada a limpar o
\VUgras{quarto}.

%%K t06-c1-02-1 p
2. --- Sobre o \VUgras{lavatório} \textsuperscript{(1)} estão espalhados os
diferentes objetos com que uma pessoa se lava, se penteia, numa palavra,
se arranja. São a \VUgras{bacia} \textsuperscript{(2)} e o \VUgras{jarro}
\textsuperscript{(3)}, a \VUgras{escova do cabelo} \textsuperscript{(4)}, o
\VUgras{pente} \textsuperscript{(5)}, o \VUgras{sabonete} \textsuperscript{(6)} e a
\VUgras{esponja} \textsuperscript{(7)}; por fim, um \VUgras{frasco}
\textsuperscript{(8)} de elixir dentífrico.

%%K t06-c1-03-1 p
3. --- No chão, diante do lavatório está o \VUgras{balde} \textsuperscript{(9)}
para recolher as águas sujas.

%%K t06-c1-04-1 p
4. --- Na \VUgras{antecâmara} \textsuperscript{(10)} vemos alguns \VUgras{cabides}
\textsuperscript{(13)} e dois \VUgras{guarda-chuvas} \textsuperscript{(11)} no
\VUgras{bengaleiro} \textsuperscript{(12)}.

%%K t06-c1-05-1 p
5. --- Do outro lado da porta está o \VUgras{guarda-fatos de espelho}
\textsuperscript{(14)}, que guarda a \VUgras{roupa branca} \textsuperscript{(15)}.

%%K t06-c1-06-1 p
6. --- Aproveitemos que a \VUgras{criada de quarto} \textsuperscript{(16)} faz a
\VUgras{cama} \textsuperscript{(17)} para estudar as suas diferentes partes. A
\VUgras{armação da cama} \textsuperscript{(20)} tem um bom \VUgras{estrado}
\textsuperscript{(21)} sobre o qual Justina vai pôr já a seguir um
\VUgras{colchão} \textsuperscript{(22)} de lã. Depois tirará das cadeiras os
dois \VUgras{lençóis} \textsuperscript{(23)} e no \VUgras{cobertor}
\textsuperscript{(24)}. A seguir colocará à cabeceira da cama o
\VUgras{travesseiro} \textsuperscript{(25)} e a \VUgras{almofada}
\textsuperscript{(26)}, que estão com o \VUgras{edredão} \textsuperscript{(27)} sobre o
\VUgras{tapete de cama} \textsuperscript{(32)}. Um espanador para sacudir o pó
das \VUgras{cortinas} \textsuperscript{(19)} e do \VUgras{dossel}
\textsuperscript{(18)}, e já está feita uma parte do trabalho.

%%K t06-c1-07-1 p
7. --- Depois terá de arrumar um pouco a \VUgras{mesa de cabeceira}
\textsuperscript{(28)}, dar corda ao \VUgras{despertador} \textsuperscript{(29)},
preparar a \VUgras{lamparina} \textsuperscript{(30)} e levar a \VUgras{vela}
\textsuperscript{(31)} para limpar a palmatória.

%%K t06-tit-2 sub
\VUsaut{5.0mm}
\VUpk{10.2pt}{40}{O cesto da costura e os acessórios da lareira.}
\VUsaut{3.0mm}

%%K t06-c1-08-1 p
8. --- O \VUgras{cesto da costura} \textsuperscript{(34)} que está junto ao
\VUgras{banco} \textsuperscript{(33)} também precisa muito de que o arrumem.
Será preciso dobrar o \VUgras{bordado} \textsuperscript{(35)}, guardar na
\VUgras{mesa de costura} \textsuperscript{(38)} o \VUgras{dedal}, a \VUgras{linha}
\textsuperscript{(36)}, as \VUgras{agulhas} \textsuperscript{(37)} e a \VUgras{tesoura}
\textsuperscript{(39)}, e fechar com a \VUgras{chave} \textsuperscript{(40)} a tampa da
mesa.

%%K t06-c1-09-1 p
9. --- Na \VUgras{mesa redonda} \textsuperscript{(41)} do centro, Justina renovará
a água da \VUgras{garrafa de cristal} \textsuperscript{(42)}. Porá
\VUgras{açúcar} no \VUgras{açucareiro} \textsuperscript{(43)}, limpará a
\VUgras{pinça do açúcar} \textsuperscript{(44)} e levará a \VUgras{escova da
roupa} \textsuperscript{(46)}.

%%K t06-c1-10-1 p
10. --- O lado direito do quarto dar-lhe-á menos trabalho, porque a
\VUgras{espreguiçadeira} \textsuperscript{(47)} e a sua \VUgras{almofada}
\textsuperscript{(48)} estão no seu sítio e, como não está frio, nada se moveu
entre os acessórios da \VUgras{lareira} \textsuperscript{(49)}. \VUgras{Pá}
\textsuperscript{(50)}, \VUgras{tenazes} \textsuperscript{(51)}, \VUgras{cães da lareira}
\textsuperscript{(55)} e \VUgras{guarda-cinzas} \textsuperscript{(56)} estão limpos; o
\VUgras{guarda-fogo} \textsuperscript{(54)} não se moveu e a \VUgras{caixa da
lenha} \textsuperscript{(52)} está bem provida de \VUgras{lenha}
\textsuperscript{(53)}.

%%K t06-noto-1 noto
\VUnotes{143.2mm}{%
(*) Babo = criança pequena; inglês \textit{baby}, francês
\textit{bébé}, italiano \textit{bambino}.%
}

%%K t06-noto-2 noto
\VUnotes{143.2mm}{%
(*) Lambrequino = francês \textit{lambrequin}, espanhol
\textit{lambrequines}, português \textit{lambrequins}, e até alemão e
inglês \textit{lambrequins}; isto é, alemão, inglês, francês, português
e espanhol por igual.%
}

%%K t06-c1-11-1 p
11. --- Terá ainda de tirar o pó do \VUgras{relógio de mesa}
\textsuperscript{(57)}, dos \VUgras{candelabros} \textsuperscript{(58)} e dos
\VUgras{esvazia-bolsos} \textsuperscript{(59)}, e depois limpar o \VUgras{espelho}
\textsuperscript{(60)}.

%%K t06-c1-12-1 p
12. --- Quanto ao \VUgras{berço} \textsuperscript{(61)} da criança (do babo (*)),
já está pronto.

%%K t06-tit-3 sub
\VUsaut{5.0mm}
\VUpk{10.2pt}{40}{A Sala de visitas.}
\VUsaut{3.0mm}

%%K t06-c2-01-1 p
1. --- Passemos agora à sala de visitas. É uma \VUgras{sala} muito bonita. A
lareira, que está no canto da direita, vai enfeitada com uma fina
\VUgras{estatueta} \textsuperscript{(1)} e dois belos \VUgras{jarrões}
\textsuperscript{(3)}, e drapeada com um \VUgras{lambrequim} \textsuperscript{(2)} de
veludo (*).

%%K t06-c2-02-1 p
2. --- Para evitar que as fagulhas do fogo peguem fogo ao tapete,
pôs-se-lhe diante um \VUgras{guarda-fogo} \textsuperscript{(4)} de cobre.

%%K t06-c2-03-1 p
3. --- Ao lado, uma elegante \VUgras{escrivaninha} \textsuperscript{(5)} de
senhora está aberta. É ali que a \VUgras{dona da casa} \textsuperscript{(23)}
escreve as suas cartas.

%%K t06-c2-04-1 p
4. --- A janela vai guarnecida com \VUgras{cortinados} \textsuperscript{(6)}
recolhidos em \VUgras{abraçadeiras} \textsuperscript{(8)} penduradas nos
\VUgras{ganchos} \textsuperscript{(7)}, e com grandes cortinas de tecido rematadas
por uma \VUgras{sanefa} \textsuperscript{(9)}.

%%K t06-c2-05-1 p
5. --- No vão da janela, uma \VUgras{floreira} \textsuperscript{(11)} contém uma
\VUgras{planta} \textsuperscript{(10)} verde, uma palmeira magnífica cujas folhas
chegam quase ao \VUgras{candeeiro de petróleo} \textsuperscript{(12)}.

%%K t06-c2-06-1 p
6. --- O \VUgras{quebra-luz} \textsuperscript{(13)} do candeeiro foi feito por
Maria, a encantadora \VUgras{jovem} \textsuperscript{(14)} que está ao
\VUgras{piano} \textsuperscript{(15)}. Sentada muito direita no \VUgras{banco}
\textsuperscript{(16)}, estuda uma peça. É muito hábil e muito cuidadosa, como o
prova o seu \VUgras{móvel da música} \textsuperscript{(17)}, que está
perfeitamente arrumado.

%%K t06-tit-4 sub
\VUsaut{5.0mm}
\VUpk{10.2pt}{40}{O Traquinas.}
\VUsaut{3.0mm}

%%K t06-c2-07-1 p
7. --- O seu irmão Paulo procura um \VUgras{livro} \textsuperscript{(19)} na
\VUgras{estante} \textsuperscript{(18)}. É um \VUgras{rapaz} \textsuperscript{(20)}
irrequieto e insuportável. Ao pendurar-se na estante para alcançar o
livro que está demasiado alto, arrisca-se a deitar abaixo o \VUgras{busto}
\textsuperscript{(21)} que está por cima.

%%K t06-c2-08-1 p
8. --- Aproveita, para fazer essa tolice, que a sua mãe está ocupada a
conversar com uma amiga, uma \VUgras{senhora} \textsuperscript{(22)} que veio
passar uns dias com ela. A mãe está sentada no \VUgras{sofá}
\textsuperscript{(24)}, ao lado da \VUgras{poltrona} \textsuperscript{(25)}, e a sua
amiga está na \VUgras{cadeira} \textsuperscript{(27)}, da qual só vemos o
\VUgras{espaldar} \textsuperscript{(26)}.

%%K t06-c2-09-1 p
9. --- A grande \VUgras{moldura} \textsuperscript{(30)} pendurada na parede contém
o \VUgras{retrato} \textsuperscript{(29)} de uma parenta. No \VUgras{álbum}
\textsuperscript{(32)} que está sobre a mesa reúnem-se as \VUgras{fotografias}
\textsuperscript{(33)} dos demais membros da família. As crianças folhearam-no
há pouco, porque as \VUgras{cadeiras} \textsuperscript{(31)} continuam agrupadas à
volta da mesa.

%%K t06-c2-10-1 p
10. --- O \VUgras{jarrão chinês} \textsuperscript{(34)} de porcelana que está
junto ao \VUgras{abre-cartas} \textsuperscript{(35)} foi oferecido a Maria pelo
seu dia de anos, e o \VUgras{biombo} \textsuperscript{(36)} que enfeita o canto da
sala foi trazido da China pelo seu tio.

%%K t06-c2-11-1 p
11. --- Alegrada pelos belos \VUgras{quadros} \textsuperscript{(37)} pendurados no
\VUgras{tabique} \textsuperscript{(42)}, iluminada pelo \VUgras{lustre}
\textsuperscript{(38)}, cujas \VUgras{lâmpadas} \textsuperscript{(39)} dão à noite uma
luz tão alegre, esta sala parece muito agradável. O fofo \VUgras{tapete}
\textsuperscript{(40)} estendido sobre o soalho e a \VUgras{campainha elétrica}
\textsuperscript{(41)}, tão prática, acabam de a tornar muito cómoda.

%%K t06-tit-5 sub
\VUsaut{3.6mm}
\VUpk{10.2pt}{40}{A Sala de jantar.}
\VUsaut{3.0mm}

%%K t06-c3-01-1 p
1. --- Soou a hora do jantar. Toda a família, salvo Maria, está reunida à
volta da mesa. A sala de jantar está agradavelmente aquecida por um
excelente \VUgras{fogão} \textsuperscript{(1)} de louça, com um \VUgras{tubo}
\textsuperscript{(2)} delgado.

%%K t06-c3-02-1 p
2. --- Esta sala não tem muitos móveis. Ao longo da parede pende, por
cima do \VUgras{escabelo} \textsuperscript{(4)}, uma \VUgras{fontezinha}
\textsuperscript{(3)} decorativa. Muito perto está o \VUgras{aparador}
\textsuperscript{(5)}, sobre o qual se deixam os pratos frios, as sobremesas e
os diferentes utensílios de mesa que de momento não fazem falta.

%%K t06-c3-03-1 p
3. --- Assim vemos, na prateleira de cima, um \VUgras{frango assado}
\textsuperscript{(7)}, um \VUgras{peixe} \textsuperscript{(8)} e uma \VUgras{taça}
\textsuperscript{(10)} com \VUgras{fruta} \textsuperscript{(9)}. Por baixo estão o
\VUgras{queijo} \textsuperscript{(11)}, o \VUgras{pão} \textsuperscript{(12)} e os
\VUgras{galheteiros} \textsuperscript{(13)}, que trazem tudo o que é preciso para
temperar a salada: o azeite, o vinagre, o \VUgras{sal} \textsuperscript{(14)} e a
\VUgras{pimenta} \textsuperscript{(15)}. Por fim, na prateleira de baixo há um
\VUgras{bolo} \textsuperscript{(17)} junto ao \VUgras{mostardeiro}
\textsuperscript{(16)}.

%%K t06-c3-04-1 p
4. --- O relógio da sala de jantar, a que se chama \VUgras{relógio de
parede} \textsuperscript{(20)}, tem uma forma muito particular. Vai encaixado
numa moldura de madeira que contém ainda um \VUgras{barómetro}
\textsuperscript{(19)} e um \VUgras{termómetro} \textsuperscript{(18)}.

%%K t06-tit-6 sub
\VUsaut{4.6mm}
\VUpk{10.2pt}{40}{O Serviço do jantar.}
\VUsaut{3.0mm}

%%K t06-c3-05-1 p
5. --- Todas as paredes da sala vão revestidas de \VUgras{painéis de
madeira} \textsuperscript{(21)} de carvalho encerado. Um \VUgras{reposteiro}
\textsuperscript{(22)} de tecido fecha bastante bem a porta que dá para a
galeria. Ali irá a família tomar o café depois de jantar. As
\VUgras{chávenas} \textsuperscript{(24)} e os \VUgras{pires} \textsuperscript{(25)} já
estão preparados no \VUgras{louceiro} \textsuperscript{(23)}.

%%K t06-c3-06-1 p
6. --- Por cima da mesa está fixado ao teto um \VUgras{candeeiro de
suspensão} \textsuperscript{(26)}, cuja luz muito viva alumia à noite toda a
sala de jantar.

%%K t06-c3-07-1 p
7. --- Ocupemo-nos agora da própria \VUgras{mesa de jantar}
\textsuperscript{(27)}. Quando a família entrou, a mesa estava posta. Sobre a
\VUgras{toalha} \textsuperscript{(28)} tinham-se colocado, no lugar de cada
comensal, um \VUgras{prato} \textsuperscript{(33)}, um \VUgras{guardanapo}
\textsuperscript{(29)}, uma \VUgras{colher} \textsuperscript{(34)}, um \VUgras{garfo}
\textsuperscript{(35)}, uma \VUgras{faca} \textsuperscript{(36)} e um \VUgras{copo}
\textsuperscript{(32)}. Havia também \VUgras{garrafas} \textsuperscript{(30)} para o
vinho e uma \VUgras{garrafa de cristal} \textsuperscript{(31)} para a água.

%%K t06-c3-08-1 p
8. --- O \VUgras{criado} \textsuperscript{(39)} trouxe primeiro a \VUgras{terrina}
\textsuperscript{(37)}, na qual ainda fumega um pouco de sopa. Agora serve o
primeiro \VUgras{prato} \textsuperscript{(38)}, um prato de carne. Depois
apresentará os que já nomeámos e, por fim, depois da \VUgras{sobremesa},
aproveitará que os seus senhores terão passado à galeria para levantar a
mesa e pôr em ordem a sala de jantar.

%%K t06-c3-08-2 p
\textit{Nota.} --- \textit{As palavras correntes relativas à comida
dão-se aqui de um modo muito sumário. A lista completar-se-á nos dois
quadros «O Hotel» (n.º 13) e «O Mercado» (n.º 15).}

%%K t06-tit-7 sub
\VUsaut{4.6mm}
\VUpk{10.2pt}{40}{A Cozinha.}
\VUsaut{3.0mm}

%%K t06-c4-01-1 p
1. --- Na cozinha, que espreitamos pela porta
entreaberta, encontramos uma desordem pitoresca. Diante do \VUgras{fogão
de chão} \textsuperscript{(1)}, onde crepita o \VUgras{fogo} \textsuperscript{(3)}, está
montado o \VUgras{assador giratório} \textsuperscript{(5)}. Um belo \VUgras{assado}
\textsuperscript{(7)} está no \VUgras{espeto} \textsuperscript{(6)} e acaba de estar
pronto.

%%K t06-c4-02-1 p
2. --- O \VUgras{fole} \textsuperscript{(8)} e a \VUgras{pá do carvão}
\textsuperscript{(9)}, que se acabaram de usar, ficaram no chão, tal como as
\VUgras{achas} \textsuperscript{(10)}.

%%K t06-c4-03-1 p
3. --- A prateleira da chaminé está arrumada: a \VUgras{lanterna}
\textsuperscript{(11)}, o \VUgras{fosforeiro} \textsuperscript{(13)}, com os
\VUgras{fósforos} \textsuperscript{(12)} dentro, o \VUgras{castiçal}
\textsuperscript{(14)} e a \VUgras{palmatória} \textsuperscript{(15)} com a sua
\VUgras{vela} \textsuperscript{(16)} estão no seu sítio. Mas o grande
\VUgras{balde do carvão} \textsuperscript{(18)}, cheio do \VUgras{carvão de pedra}
\textsuperscript{(17)} que a \VUgras{cozinheira} \textsuperscript{(23)} acaba de ir
buscar à adega, ficou no meio da cozinha.

%%K t06-c4-04-1 p
4. --- A verdade é que a pobre mulher tem de se despachar para que a
comida esteja pronta a tempo. Agora está junto ao \VUgras{fogão}
\textsuperscript{(19)}, a mexer um \VUgras{molho} \textsuperscript{(24)} que não deve
deixar queimar.

%%K t06-c4-05-1 p
5. --- No \VUgras{forno} \textsuperscript{(20)} coze um bolo que preparou para o
lanche das crianças. Por cima do fogão pendem a \VUgras{concha}
\textsuperscript{(21)} e a \VUgras{escumadeira} \textsuperscript{(22)}.

%%K t06-tit-8 sub
\VUsaut{4.0mm}
\VUpk{10.2pt}{40}{A Bateria de cozinha.}
\VUsaut{3.0mm}

%%K t06-c4-06-1 p
6. --- A \VUgras{bateria de cozinha} \textsuperscript{(25)}, pendurada ao fundo da
sala, está muito limpa. Os belos \VUgras{tachos} \textsuperscript{(26)} de cobre, o
\VUgras{leiteiro} \textsuperscript{(27)}, o \VUgras{coador} \textsuperscript{(28)} e o
\VUgras{ralador} \textsuperscript{(29)} foram limpos com esmero.

%%K t06-c4-07-1 p
7. --- O \VUgras{saleiro} \textsuperscript{(30)} está posto sobre o pequeno
aparador, ao lado do \VUgras{bule} \textsuperscript{(31)} e do \VUgras{garrafão do
azeite} \textsuperscript{(32)}. A \VUgras{gaveta} \textsuperscript{(33)} guarda de certo
os talheres e as facas de cozinha.

%%K t06-c4-08-1 p
8. --- A \VUgras{vassoura} \textsuperscript{(34)} e o \VUgras{espanador}
\textsuperscript{(35)} estão a um canto. A cozinheira acabou de usar o
\VUgras{cepo} \textsuperscript{(36)}; deixou-o junto à janela.

%%K t06-c4-09-1 p
9. --- Uma \VUgras{toalha de mãos} \textsuperscript{(37)} pende da parede, ao lado
do \VUgras{funil} \textsuperscript{(38)}. Nesta parte da cozinha guardam-se a
\VUgras{grelha} \textsuperscript{(39)}, o \VUgras{cutelo de picar}
\textsuperscript{(40)} e a \VUgras{tábua de picar} \textsuperscript{(41)}. Depois, por
cima do cutelo, numa \VUgras{prateleira} \textsuperscript{(42)}, há
\VUgras{panelas} \textsuperscript{(43)}, a \VUgras{caçarola} \textsuperscript{(44)} com a
sua \VUgras{tampa} \textsuperscript{(45)} e a \VUgras{marmita} \textsuperscript{(46)}. A
\VUgras{frigideira} \textsuperscript{(47)} pende ao lado.

%%K t06-c4-10-1 p
10. --- Só a \VUgras{torneira} \textsuperscript{(48)} de cobre põe um brilho neste
canto. O \VUgras{lava-loiça} \textsuperscript{(49)}, sobre o qual está o grande
\VUgras{cântaro} \textsuperscript{(50)}, deve ser muito cómodo com o seu
armarinho, no qual se guardam o \VUgras{tonel do vinagre}
\textsuperscript{(51)} com a sua \VUgras{torneirinha} \textsuperscript{(52)}, a
\VUgras{caixa do lixo} \textsuperscript{(53)} e a \VUgras{loiça suja}
\textsuperscript{(54)}.

%%K t06-tit-9 sub
\VUsaut{4.0mm}
\VUpk{10.2pt}{40}{Outras coisas que se guardam na cozinha.}
\VUsaut{3.0mm}

%%K t06-c4-11-1 p
11. --- O grande \VUgras{alguidar} \textsuperscript{(55)} no qual se lavam os
\VUgras{legumes} \textsuperscript{(58)} e o \VUgras{caldeirão} \textsuperscript{(56)} de
ferro fundido posto sobre o \VUgras{ladrilhado} \textsuperscript{(57)} devem também
ter o seu lugar nesse armário.

%%K t06-c4-12-1 p
12. --- À cozinheira não faltam fogareiros, porque há ainda, na esquina
da mesa, um \VUgras{fogareiro a gás} \textsuperscript{(59)} no qual se aquece água
num tacho pequeno. O \VUgras{vapor} \textsuperscript{(60)} que dele sai indica-nos
que deve estar a ferver. De certo essa água servirá para o café, porque
na outra ponta da mesa estão a \VUgras{cafeteira} \textsuperscript{(64)} e o
\VUgras{moinho de café} \textsuperscript{(63)}.

%%K t06-c4-13-1 p
13. --- O fogareiro está ligado ao \VUgras{bico de gás} \textsuperscript{(62)} por
um longo \VUgras{tubo de borracha} \textsuperscript{(61)}.

%%K t06-c4-14-1 p
14. --- A \VUgras{mulher-a-dias} \textsuperscript{(66)}, que vem ajudar a
cozinheira, prepara os legumes para o jantar. Quando estiverem
descascados e lavados, pô-los-á na \VUgras{caçoila} \textsuperscript{(65)} até à
hora de os cozer.

%%K t06-tit-10 sub
\VUsaut{4.0mm}
{\centering\fontsize{10.2pt}{10.2pt}\selectfont\textls[40]{\scshape A Casa de Banho.} \textit{(A sala do banho.)}\par}
\VUsaut{3.0mm}

%%K t06-c5-01-1 p
1. --- Só nos resta uma indiscrição por cometer para conhecer os
principais compartimentos da casa: ver a instalação da casa de banho.
Não será longo, porque não é grande, e saberemos depressa que a
\VUgras{banheira} \textsuperscript{(1)} tem um bom \VUgras{esquentador}
\textsuperscript{(2)}, que a \VUgras{saboneteira} \textsuperscript{(3)} está ao alcance
da mão de quem se banha e que o \VUgras{toalheiro} \textsuperscript{(5)} permite
secar muito facilmente as \VUgras{toalhas} \textsuperscript{(4)}.

%%K t06-c5-02-1 p
2. --- Esta sala tem os tabiques cobertos de \VUgras{azulejos}
\textsuperscript{(6)}, o que permitiu instalar um \VUgras{duche}
\textsuperscript{(7)} sem recear que a água que salpica ao usá-lo estrague as
paredes. Uma pequena \VUgras{bacia} \textsuperscript{(8)} de zinco, que serve para
os banhos de pés, completa o mobiliário.
\VUsaut{6.0mm}
\VUfilet{22mm}
